\documentclass{article}
\usepackage[T1]{fontenc}
\usepackage{lmodern}
\usepackage{graphicx}
\usepackage{amsmath,amssymb}
\usepackage[a4paper,margin=2cm]{geometry}
\usepackage[absolute,overlay]{textpos}
\setlength{\TPHorizModule}{1mm}
\setlength{\TPVertModule}{1mm}
\usepackage[indonesian]{babel}
\usepackage{hyperref}
\setlength{\emergencystretch}{2em}
% unit: Halaman 20

\begin{document}

% === Halaman 20 (python/native) ===

20

Contoh parameter RSA 

• Modulus n sepanjang 1024 bit (setara 300 angka decimal

• Bilangan prima p dan q masing-masing panjangnya sekitar  154 angka decimal

• Sumber:  https://www.di-mgt.com.au/rsa\_alg.html

• n = 

1192941348401695090555272113312556496446065696615276380120674819549430568

5115033380631595703771562029730500011862877084668996911289221224545711806

0574995989517080042105263427376322274266393116193517839570773505632231596

6811219273374739732203125125990612313222509455062600665575382385175753906

21262940383913963

• p = 

1093376618363257581761151703473066828715579998463222345413874567112127345

6287670008290843302875521274970245314593222946129064538358581018615539828

479146469

• q = 

1091061696734911023172373407861492264533706088214174896820983422513897601

1179993394299810159736904468554021708289824396553412180514827996444845438

176099727

\end{document}
